\section{The Judgment First}

On 19 September 1851, the fifth anniversary of the reported
apparition, Philibert de Bruillard, Bishop of Grenoble, closed five
years of diocesan inquiry with a \term{mandement} bearing a
doctrinal judgment on the event of La Salette. Everything in this
study hangs on reading its operative sentence exactly---and on
noticing how narrow it is.

\begin{operativetext}{Mandement of the Bishop of Grenoble, 19
September 1851, article~1, in the witnesses this study could reach}
\textbf{French, as transcribed by the Sanctuary of Notre-Dame de La
Salette (the ellipsis is the Sanctuary's own):}\par
«~L'apparition de la Sainte Vierge à deux bergers sur la montagne
de La Salette [\ldots] porte en elle-même tous les caractères de la
vérité et que les fidèles sont fondés à la croire indubitable et
certaine.~»\par
\medskip
\textbf{The same operative clause as quoted by the Dicastery for the
Doctrine of the Faith in the presentation of its 2024 norms:}\par
«~Les fidèles sont fondés à la croire indubitable et certaine~»
(``Decree of the Bishop of Grenoble, 19 September 1851'').\par
\medskip
\textbf{Article~1 entire, in the English of an 1854 witness:}\par
``We judge that the Apparition of the Blessed Virgin to two
cowherds, on the 19th of September, 1846, on a mountain of the chain
of the Alps, situated in the parish of La Salette, in the
arch-presbytery of Corps, bears within itself all the
characteristics of truth, and that the faithful have grounds for
believing it indubitable and
certain.''\endnote{French: Sanctuaire Notre-Dame de La Salette,
``L'histoire,'' \url{https://lasalette.cef.fr/lhistoire/}, read
2026-07-25; the bracketed ellipsis is printed by the Sanctuary and
is reproduced rather than filled. Holy See witness to the same
clause: Dicastery for the Doctrine of the Faith, \work{Norms for
Proceeding in the Discernment of Alleged Supernatural Phenomena}
(17 May 2024), presentation, ``Reasons for the New Norms,'' English
text at
\url{https://www.vatican.va/roman_curia/congregations/cfaith/documents/rc_ddf_doc_20240517_norme-fenomeni-soprannaturali_en.html}
(read 2026-07-25), which prints the French clause and identifies it
as the ``Decree of the Bishop of Grenoble, 19 September 1851.''
English of the whole article: William Bernard Ullathorne, \work{The
Holy Mountain of La Salette: A Pilgrimage of the Year 1854} (London:
Richardson, 1854), pp.~120--121, read 2026-07-25 in the Internet
Archive scan (References). A French printing of the integral
mandement was not located by this project; the French clause is
therefore given only as far as two independent witnesses print it,
and no French reconstruction of the whole article is offered
(Appendix~A). An 1866 Grenoble life of the diocesan
\term{rapporteur} reports the judgment in the same French
words---``porte en elle-même tous les caractères de la vérité, et que les
fidèles sont fondés à la croire indubitable et certaine''---inside
the biographer's own sentence rather than as a transcription of the
act. All web loci in these notes were read on 2026-07-25 unless
another date is given. English renderings of French, Latin, and
Italian are editorial working translations, visibly commentary, and
never received translation witnesses, except where an English text
of the Holy See or of an identified nineteenth-century translator is
expressly quoted as such.}
\end{operativetext}

The act did not stop at article~1. An 1854 Italian translation of
the pilgrim's manual written by the diocesan \term{rapporteur}
prints eight articles: article~1 as above; article~2, that the fact
gains a further degree of certainty from the immense and spontaneous
concourse of the faithful and from the multitude of prodigies that
followed; article~3, authorizing the cult of Our Lady of La Salette
and permitting it to be preached with its practical and moral
consequences; article~4, forbidding the publication of any
particular formula of prayer, any hymn, any book of devotion without
prior written episcopal approval; article~5, expressly forbidding
the faithful and the priests of the diocese ever to rise publicly,
in speech or in writing, against the fact now published; article~6,
announcing the purchase of the ground and the intention to build
there a church and a hospice for pilgrims, with an appeal for alms
and a commission of priests and laymen to oversee the building and
the offerings; article~7, the pastoral exhortation to penance, to
divine worship, and to the sanctification of Sunday, which the act
calls the apparition's principal end; and article~8, ordering the
mandement read and published in every church and chapel of the
diocese at the parish or community Mass on the Sunday following its
receipt. It is dated at Grenoble under the bishop's signature, the
seal of his arms, and his secretary's countersignature, ``19
September 1851 (fifth anniversary of the celebrated
apparition).''\endnote{Pierre-Joseph Rousselot, \work{Manuale del
pellegrino a Nostra Signora di La Salette} (Italian translation of
the \work{Manuel du pèlerin}), 1854, pp.~42--47, read 2026-07-25 in
the Internet Archive scan of the Bavarian State Library copy
(References). The articles are summarized here, not translated
verbatim from the Italian; the English of arts.~1 and 3--4 and 7 in
Ullathorne 1854, pp.~120--121, agrees in substance with this Italian
witness, and the two were compared. Article~5's prohibition of
public opposition is in the Italian witness; Ullathorne's abridged
English renders the same effect as a prohibition on publicly
contradicting or calling the fact in question. Where the two
witnesses differ in fullness this study follows neither silently:
the article inventory is reported as the Italian witness prints it
and the difference is recorded in Appendix~A.}

That is the whole positive judgment: a declaration that the
apparition to two shepherd children on a named mountain on a named
day bears in itself all the marks of truth, that the faithful have
grounds for believing it indubitable and certain, an authorization
of cult, a reservation of devotional publications to written
episcopal approval, a disciplinary prohibition of public contest
within the diocese, and a building programme. The act is diocesan,
made by the local ordinary under the discipline of his own century.
It is not a papal definition, not an exercise of infallibility, not
a source of new public Revelation, and---this is the point on which
the whole later history turns---not an approval of any secret, any
later publication, or any prediction. Twenty-six years later the
Sacred Congregation of Rites said in so many words what Rome
considered the status of La Salette and Lourdes alike to be: such
apparitions or revelations have been ``neither approved nor
reprobated or condemned by the Apostolic See, but only permitted as
to be piously believed by merely human faith'' (\S10.2).

\studythesis{This is an ecclesial-judgment-first study, not a
devotional narrative and not an exposé. It reconstructs the reported
encounter of 19 September 1846 from the nearest witnesses this
project could actually reach; it separates the public discourse as
first recorded from the received forms in which it now circulates,
and shows what that separation does and does not license; it quotes
the 1851 mandement at its operative clause and states its exact
object; it treats the secrets as a juridical and documentary
problem---two short manuscripts written for Pius IX in July 1851,
then Mélanie Calvat's progressively expanded publications
culminating in the Lecce booklet of 1879 printed under a diocesan
\term{imprimatur}, then the Holy Office decree published at
\work{AAS} 7 (1915) 594 forbidding all the faithful to treat of the
subject, then the condemnation of a 1922 reprint at \work{AAS} 15
(1923) 287--288. It then publishes the 1879 text and analyzes it,
under the disclosure set out in full at \S5.1: the 1915 decree is
quoted at its command, its penalties, and its saving clause, its
present force is reported as unresolved rather than decided in
either direction, and the analysis is published on the project
maintainer's express direction as the editorial act of a private
study and not as any ecclesiastical judgment. The question that
analysis answers is whether the text published in 1879 is the text
sent to Pius~IX in 1851; the answer reached, labelled throughout as
project synthesis, is that it is not (\S6.9). Public Revelation is complete in Christ; no La
Salette event adds to it, and an approved private revelation asks
only prudent human assent (\term{fides humana}), never the divine
and Catholic faith owed to Revelation. A printing permission is not
an authentication; a Roman archive's receipt is not papal approval;
a religious institute, a basilica, a papal letter, and an optional
memorial each establish exactly themselves. Where the record is
thinner than the received narrative---above all in the absence of an
accessible French printing of the integral mandement and in the
contested reliability that attaches to both seers' later years---the
ceiling is stated rather than papered over.}

The study runs in the order the evidence demands: the event and its
earliest transmission (\S2); the public message and its two textual
strata (\S3); the inquiry and the 1851 judgment (\S4); the secrets,
their custody, their publication, and the Roman acts of 1880, 1915
and 1923 with the juridic question they leave open (\S5); the 1879
text published and analyzed against the 1851 manuscripts, with the
judgment this study is willing to defend (\S6); the seers'
later lives and the contested reliability that attaches to them
(\S7); the shrine, the Missionaries of 1852, and the discipline of
cult (\S8); papal and liturgical reception (\S9); the doctrinal
boundary---what approval does and does not oblige (\S10); and the
current standing of the whole dossier under the 2024 norms (\S11).
Work-wide bounds, taxonomy, and qualifications are gathered in the
terminal appendix; witnesses and rights are in the References.
